\documentclass[12pt,a4paper]{article}

\usepackage{iftex}
\ifPDFTeX
  \usepackage[T2A]{fontenc}
  \usepackage[utf8]{inputenc}
  \usepackage[russian,english]{babel}
\else
  \usepackage{fontspec}
  \IfFontExistsTF{Times New Roman}
    {\setmainfont{Times New Roman}}
    {\setmainfont{TeX Gyre Termes}}
  \IfFontExistsTF{Consolas}
    {
      \setmonofont{Consolas}
      \newfontfamily\cyrillicfonttt{Consolas}
    }
    {
      \setmonofont{Courier New}
      \newfontfamily\cyrillicfonttt{Courier New}
    }
  \usepackage{polyglossia}
  \setdefaultlanguage{russian}
  \setotherlanguage{english}
\fi

\usepackage{geometry}
\geometry{left=25mm,right=20mm,top=20mm,bottom=20mm}
\usepackage{array}
\usepackage{booktabs}
\usepackage{tabularx}
\usepackage{hyperref}
\usepackage{xurl}

\newcolumntype{Y}{>{\raggedright\arraybackslash}X}
\newcommand{\gap}{$\to$}
\emergencystretch=2em

\title{Автоматизация итеративного ревью кода с помощью LLM-комментариев для снижения числа итераций исправления}
\author{А. К. Белоногов \and А. А. Лаушкина \and В. Д. Волоха \and И. З. Смирнов}
\date{2026}

\begin{document}

\maketitle

\begin{abstract}
В работе рассматривается задача автоматического ревью кода внутри итеративного
исправления программ. В таком цикле агент-разработчик получает требование,
делает первую правку, проверяет ее тестами и при необходимости продолжает
исправление. Основная гипотеза состоит в том, что один содержательный
LLM-комментарий к первой правке может уменьшить число дальнейших итераций или
предотвратить ухудшение результата. Для проверки этой гипотезы реализован
прототип MergeMind: генератор формирует возможное замечание, оценщик решает,
достаточно ли оно связано с видимым изменением, а копирайтер переводит принятое
замечание в короткую инструкцию для агента. SWE-CI используется как
экспериментальный стенд, потому что он измеряет не сходство текста комментария с
историческим ревью, а последующее поведение агента в цикле тестирования.
Эксперименты дают смешанную картину: на задаче \texttt{cle-b/httpdbg}
MergeMind снизил финальный разрыв по тестам с 11 до 5, но на
\texttt{inline-snapshot} не уменьшил финальный разрыв и число итераций. На
четырехзадачном проверочном запуске устойчивого улучшения также не получено.
Следовательно, автоматический комментарий может помогать отдельным
исправлениям, но гипотеза о стабильном снижении числа итераций требует более
широкой проверки.
\end{abstract}

\section{Введение}

Ревью кода является важным способом контроля качества программного обеспечения:
оно помогает находить дефекты, регрессии, неполные исправления и несоответствие
изменения исходному требованию. Современные модели уже могут автоматически
формировать замечания по изменению кода, однако практические исследования
показывают, что такие замечания часто бывают шумными: они могут быть слишком
общими, плохо связанными с дефектом или непригодными для следующей
правки~\cite{acrpractice,desiview}.
Однако практическая ценность такого комментария определяется не только тем,
похож ли он на человеческий текст. Комментарий может выглядеть убедительно, но
не помогать исправлению, указывать на несущественную проблему или направлять
агента к вредной дополнительной правке. Поэтому в работах о желательных
комментариях ревью отдельно подчеркивается необходимость отбирать полезные и
обоснованные замечания, а не просто увеличивать их количество~\cite{desiview}.

В этой статье автоматический комментарий рассматривается не как финальный ответ
для человека, а как средство управления итеративным исправлением кода. Такой
сценарий отличается от обычной генерации ревью: важно не только качество текста,
но и то, меняется ли поведение агента после получения комментария. Если
замечание действительно полезно, агент должен быстрее прийти к рабочему
варианту, удержать меньше падающих тестов или хотя бы не ухудшить уже найденное
частичное решение.

Проверять такую гипотезу удобно в среде, где каждая попытка исправления
измеряется тестами. В работе используется SWE-CI~\cite{sweci}: агент получает
требование к изменению, меняет код, а результат оценивается запуском тестов в
цикле непрерывной интеграции. В отличие от статических наборов задач на
исправление кода, такой стенд позволяет наблюдать траекторию решения: сколько
итераций выполнил агент, как менялось число падающих тестов и не появлялись ли
новые регрессии.

Исследовательский вопрос работы формулируется следующим образом: может ли один
автоматически сформированный комментарий ревью, переданный агенту после первой
правки, улучшить дальнейший ход исправления? Из него выделяются три проверяемые
гипотезы:

\begin{itemize}
  \item \textbf{H1}: комментарий уменьшает число итераций, необходимых агенту
  для достижения лучшего наблюдаемого результата;
  \item \textbf{H2}: комментарий уменьшает финальное или лучшее число падающих
  тестов по сравнению с запуском без MergeMind;
  \item \textbf{H3}: комментарий не должен давать улучшение только ценой
  неконтролируемых дополнительных правок и существенного роста стоимости
  запуска.
\end{itemize}

Вклад работы состоит в реализации и первичной проверке цепочки из генератора,
оценщика и копирайтера для одного итогового комментария внутри цикла SWE-CI.
Результаты не трактуются как доказательство общего превосходства метода.
Напротив, отрицательные и смешанные запуски используются для уточнения границ
подхода.

\section{Связанные работы}

\textbf{Автоматическое ревью кода.}
Работы о желательных комментариях показывают, что не каждое историческое
замечание одинаково полезно для обучения и оценки моделей~\cite{desiview}.
Практические исследования автоматического ревью также указывают на шум: внешне
похожий на человеческое ревью текст может быть слишком общим или плохо связанным
с реальным дефектом~\cite{acrpractice}. Поэтому MergeMind не стремится
сгенерировать как можно больше замечаний. Метод выбирает один комментарий,
достаточно связанный с видимым изменением и пригодный для следующей правки
агента.

\textbf{Агенты для исправления программ.}
SWE-bench стал одним из основных наборов задач для оценки LLM-агентов на
исправлении реальных ошибок в репозиториях~\cite{swebench}. SWE-agent
показывает, что агенту важны не только модель и промпт, но и организация
взаимодействия с окружением разработки~\cite{sweagent}. Agentless демонстрирует
другую линию: часть задач можно решать без сложного агентного цикла, если
аккуратно построить локализацию и исправление~\cite{agentless}. Reflexion
исследует самокоррекцию через текстовую обратную связь после неудачных
попыток~\cite{reflexion}. MergeMind связан с этим направлением, но использует
не свободное саморефлексивное описание агента, а внешний комментарий ревью по
конкретной первой правке.

\textbf{Оценка через поведение, а не только через текст.}
Для ревью кода текстовое сходство с эталонным комментарием не является
достаточным критерием: два разных комментария могут быть полезными, а похожий
комментарий может не помогать исправлению. Поэтому новые работы по
defect-focused review и CodeReviewQA смещают внимание к пониманию дефекта,
действию комментария и проверяемости результата~\cite{defectfocused,codereviewqa}.
CoDocBench дополнительно показывает важность контекста изменения для согласования
кода и сопутствующих артефактов~\cite{codocbench}. METAMON близок по идее
проверки через спецификации и тесты: модельные ответы нужно подтверждать
внешним проверяющим механизмом~\cite{metamon}. В этой работе таким механизмом
выступает SWE-CI, а Tree-sitter рассматривается как вспомогательный инструмент
структурного разбора измененных файлов~\cite{treesitter}.

\section{Постановка задачи}

Рассматривается задача исправления кода в репозитории. На вход агенту подается
требование \(r\), текущее состояние репозитория \(s_0\) и доступный контекст
проекта. Агент делает первую правку \(p_1\), после чего получается новое
состояние \(s_1\). В обычном запуске SWE-CI это состояние проверяется тестами,
и агент продолжает работу в следующей итерации, если задача еще не решена.

В запуске с MergeMind между первой правкой и запуском тестов добавляется один
дополнительный шаг. MergeMind анализирует требование \(r\), видимое изменение
\(p_1\), доступный контекст репозитория и, если уже есть наблюдаемые сведения о
падающих тестах, их краткое описание. После этого система либо возвращает один
итоговый комментарий \(c_1\), либо отказывается от комментария. Если комментарий
принят, агент получает его как инструкцию для дополнительного исправления уже
измененных файлов. После этой дополнительной правки запускаются тесты.

В этой работе \textbf{итерацией} называется один цикл работы агента между двумя
запусками тестов SWE-CI: получение требования и контекста, изменение кода,
необязательное получение комментария MergeMind, дополнительная правка по
комментарию и последующий запуск тестов. Такое определение важно, потому что
комментарий оценивается не сам по себе, а по тому, как меняется измеряемый
результат следующего запуска.

Комментарий считается допустимым только при выполнении трех условий:

\begin{itemize}
  \item он основан на видимом требовании, изменении кода и доступном контексте;
  \item он указывает на конкретную возможную ошибку в поведении программы;
  \item он пригоден для ограниченного исправления уже измененных файлов и не
  использует скрытое правильное решение.
\end{itemize}

\section{Метод}

MergeMind строится как цепочка из трех этапов:

\begin{itemize}
  \item \textbf{генератор} формирует черновой комментарий по изменению;
  \item \textbf{оценщик} проверяет, достаточно ли комментарий обоснован и
  связан с видимой ошибкой;
  \item \textbf{копирайтер} переводит принятый комментарий в короткую
  инструкцию для исправления агентом-разработчиком.
\end{itemize}

На выходе цепочки агент получает не список замечаний, а не более одного
итогового комментария за итерацию. Это сделано намеренно: цель эксперимента не
в том, чтобы нагрузить агента большим количеством текста, а в том, чтобы
проверить действие одного выбранного сигнала.

\subsection{Встраивание в цикл SWE-CI}

Последовательность запуска выглядит следующим образом:

\begin{enumerate}
  \item SWE-CI формирует требование для текущей задачи.
  \item Агент-разработчик делает первую правку кода.
  \item MergeMind получает требование, видимое изменение и доступный контекст.
  \item Генератор, оценщик и копирайтер формируют один итоговый комментарий или
  отказываются от комментария.
  \item Агент получает комментарий и список файлов, которые разрешено менять.
  \item Агент делает дополнительную правку, после чего запускаются тесты.
\end{enumerate}

Важное ограничение эксперимента: MergeMind не получает скрытое эталонное
исправление, не использует \texttt{target\_sha} и не видит скрытую разницу с
правильным решением. Система работает только с тем, что было бы доступно в
реальном цикле исправления: требованием, уже сделанной правкой, контекстом
репозитория и наблюдаемыми результатами предыдущих запусков.

Технически результат работы MergeMind сохраняется в файл с комментарием для
агента. Если оценщик не принимает замечание, дополнительная правка не
запускается. Если замечание принято, агенту передается краткая инструкция и
ограничение на файлы, которые можно менять. Такое ограничение нужно, чтобы
комментарий не превращался в свободную попытку переписать весь репозиторий.

\section{Экспериментальная схема}

SWE-CI используется как стенд, а не как тема работы. Для каждой задачи
сравниваются два режима:

\begin{itemize}
  \item \textbf{Без MergeMind}: агент работает в обычном цикле SWE-CI.
  \item \textbf{С MergeMind}: после первой правки агент получает один
  автоматически выбранный комментарий и может сделать дополнительное исправление
  перед запуском тестов.
\end{itemize}

Во всех сравнимых запусках используются одинаковые задачи, одинаковый лимит
итераций и одна локальная модель \texttt{qwen3.6-27b@iq2\_xxs} через LM Studio.
Название модели приводится для воспроизводимости: оно означает конкретную
локальную сборку LLM, которая использовалась и для агента, и для MergeMind в
данных запусках.

\subsection{Метрики}

Основные метрики выбираются так, чтобы отвечать на гипотезы H1--H3.

\begin{itemize}
  \item \textbf{Число итераций} показывает, сколько циклов исправления выполнил
  агент до завершения запуска или до достижения заданного лимита.
  \item \textbf{Итерация лучшего результата} показывает, на какой итерации
  впервые достигнут минимальный наблюдаемый разрыв по тестам.
  \item \textbf{Финальный разрыв} показывает число падающих тестов в последней
  итерации.
  \item \textbf{Лучший разрыв} показывает минимальное число падающих тестов за
  весь запуск.
  \item \textbf{EvoScore} используется как официальный показатель SWE-CI для
  сравнения хода решения.
  \item \textbf{Исправленные и новые падения} показывают, какие финальные
  падения базового запуска исчезли в запуске с MergeMind и какие новые падения
  появились.
  \item \textbf{Стоимость комментария} учитывает дополнительные вызовы модели и
  токены, затраченные на работу MergeMind.
\end{itemize}

Значение \texttt{final\_gap=-1} не считается улучшением. Оно означает, что
финальный запуск тестов нельзя интерпретировать как обычное число падающих
тестов.

\section{Результаты}

\subsection{Положительные и смешанные пробные запуски}

В таблице~\ref{tab:positive} приведены два пробных сравнения. В строке
``С MergeMind'' агент после первой правки получил один комментарий,
сформированный генератором, оценщиком и копирайтером.

\begin{table}[h]
\centering
\small
\caption{SWE-CI: сравнение запуска без MergeMind и с MergeMind}
\label{tab:positive}
\begin{tabularx}{\textwidth}{@{}Y Y c c c Y@{}}
\toprule
Задача & Режим & Итерации & Финальный разрыв & Лучший разрыв & Ход разрыва по тестам \\
\midrule
\texttt{inline-snapshot} & Без MergeMind & 3 & 12 & 12 & 16 \gap 106 \gap 13 \gap 12 \\
\texttt{inline-snapshot} & С MergeMind & 3 & 12 & 12 & 16 \gap 13 \gap 13 \gap 12 \\
\addlinespace
\texttt{cle-b/httpdbg} & Без MergeMind & 3 & 11 & 5 & 5 \gap -1 \gap 5 \gap 11 \\
\texttt{cle-b/httpdbg} & С MergeMind & 3 & 5 & 5 & 5 \gap 5 \gap 5 \gap 5 \\
\bottomrule
\end{tabularx}
\end{table}

На задаче \texttt{inline-snapshot} MergeMind не уменьшил число итераций и не
улучшил финальный разрыв: оба режима завершились с финальным разрывом 12.
При этом ход запуска стал стабильнее: без MergeMind после первой правки разрыв
вырос с 16 до 106, а с MergeMind такого резкого ухудшения не было. Официальный
EvoScore на этом пробном сравнении вырос с 0.5045 до 0.8333, но этот результат
нельзя трактовать как доказательство сокращения числа итераций.

На задаче \texttt{cle-b/httpdbg} сигнал сильнее. Базовый запуск завершился с
финальным разрывом 11, а запуск с MergeMind удержал разрыв 5 на всех
наблюдаемых итерациях. В отдельном сравнении финальных наборов падений
MergeMind устранил 6 финальных падений базового режима и не добавил новых.
Официальный EvoScore для базового запуска равен -0.5439, а для соответствующего
запуска с MergeMind равен 0.0000. Однако число итераций и здесь не снизилось:
оба завершенных запуска использовали заданный лимит из трех итераций.

\subsection{Проверка на нескольких задачах}

Чтобы проверить, не является ли положительный результат на
\texttt{cle-b/httpdbg} единичным, был проведен четырехзадачный запуск с более
строгой проверкой комментариев. Его результат оказался отрицательным:
среднее число итераций осталось 3.000 против 3.000, средний лучший разрыв
остался 6.250 против 6.250, а все финальные запуски в обеих группах были
невалидными. Среднее изменение официального EvoScore составило -0.070.

Этот запуск важен для интерпретации метода. Он показывает, что даже более
строгий оценщик комментариев не гарантирует улучшения поведения агента. Кроме
того, MergeMind добавил стоимость запуска: в данном батче на комментарии было
потрачено 135610 токенов и 32 дополнительных вызова модели. Поэтому стоимость
комментария должна учитываться вместе с качеством решения, а не отдельно от
него.

\section{Обсуждение}

Полученные результаты удобно рассматривать по исходным гипотезам.

\textbf{H1 не подтверждена.} В завершенных сравнимых запусках агент не стал
использовать меньше итераций. На \texttt{inline-snapshot} MergeMind раньше
удержал лучший разрыв, но итоговое число итераций осталось тем же. На
\texttt{cle-b/httpdbg} улучшился финальный разрыв, однако запуск также
использовал весь заданный лимит итераций.

\textbf{H2 получила частичное подтверждение.} На \texttt{cle-b/httpdbg}
комментарий помог избежать ухудшения к финалу: финальный разрыв снизился с 11
до 5, а новых финальных падений не появилось. На \texttt{inline-snapshot}
финальный и лучший разрывы не улучшились, хотя траектория стала менее резкой.
На четырехзадачном запуске улучшения лучшего разрыва не было.

\textbf{H3 остается ограничением метода.} Комментарий полезен только тогда,
когда он приводит к измеримому улучшению или предотвращает ухудшение. Если
оценщик пропускает слабое замечание, система может потратить дополнительные
вызовы модели без выигрыша по тестам. Поэтому MergeMind должен оцениваться не
как генератор красивого текста, а как часть цикла исправления, где важны
финальный результат, траектория тестов и стоимость.

Главное ограничение текущей работы -- малый размер выборки SWE-CI. Текущие
данные показывают реализуемость подхода и один сильный положительный пример, но
не дают статистически устойчивого вывода о снижении числа итераций. Следующий
необходимый шаг -- парный прогон на 20--30 задачах с одинаковыми лимитами,
одинаковой моделью и заранее зафиксированными метриками H1--H3.

\section{Заключение}

В работе представлен прототип MergeMind для автоматического комментария ревью
внутри итеративного исправления кода. Метод использует генератор, оценщик и
копирайтер, а агент получает не более одного комментария за итерацию. SWE-CI
применяется как экспериментальный стенд: он позволяет измерять
не только текст комментария, но и последующее поведение агента.

Первичные результаты являются смешанными. На одной задаче MergeMind снизил
финальный разрыв по тестам с 11 до 5 и не добавил новых финальных падений. На
другой задаче он стабилизировал ход запуска, но не улучшил финальный результат.
На четырехзадачной проверке устойчивого улучшения не получено. Поэтому текущий
вывод формулируется осторожно: LLM-комментарий может быть полезным управляющим
сигналом для отдельного исправления, но снижение числа итераций пока не
доказано и требует более широкой экспериментальной проверки.

\begin{thebibliography}{10}

\bibitem{desiview}
Y. Yu, L. Zhang, G. Rong, H. Shen, J. Zhang, H. Yan, G. Shi,
D. Shao, R. Pan, Y. Li, Q. Wang, Z. Tian.
\newblock Distilling Desired Comments for Enhanced Code Review with Large Language Models.
\newblock arXiv:2412.20340, 2024.
\newblock \url{https://arxiv.org/abs/2412.20340}

\bibitem{acrpractice}
U. Cihan, V. Haratian, A. Icoz, M. K. Gul, O. Devran,
E. F. Bayendur, B. M. Ucar, E. Tuzun.
\newblock Automated Code Review In Practice.
\newblock arXiv:2412.18531, 2024.
\newblock \url{https://arxiv.org/abs/2412.18531}

\bibitem{defectfocused}
J. Lu, L. Jiang, X. Li, J. Fang, F. Zhang, L. Yang, C. Zuo.
\newblock Towards Practical Defect-Focused Automated Code Review.
\newblock arXiv:2505.17928, 2025.
\newblock \url{https://arxiv.org/abs/2505.17928}

\bibitem{codereviewqa}
H. Y. Lin, C. Liu, H. Gao, P. Thongtanunam, C. Treude.
\newblock CodeReviewQA: The Code Review Comprehension Assessment for Large Language Models.
\newblock arXiv:2503.16167, 2025.
\newblock \url{https://arxiv.org/abs/2503.16167}

\bibitem{codocbench}
K. Pai, P. Devanbu, T. Ahmed.
\newblock CoDocBench: A Dataset for Code-Documentation Alignment in Software Maintenance.
\newblock arXiv:2502.00519, 2025.
\newblock \url{https://arxiv.org/abs/2502.00519}

\bibitem{swebench}
C. Jimenez, J. Yang, A. Wettig, S. Yao, K. Pei, O. Press, K. Narasimhan.
\newblock SWE-bench: Can Language Models Resolve Real-World GitHub Issues?
\newblock arXiv:2310.06770, 2023.
\newblock \url{https://arxiv.org/abs/2310.06770}

\bibitem{sweagent}
J. Yang, C. E. Jimenez, A. Wettig, K. Lieret, S. Yao, K. Narasimhan, O. Press.
\newblock SWE-agent: Agent-Computer Interfaces Enable Automated Software Engineering.
\newblock arXiv:2405.15793, 2024.
\newblock \url{https://arxiv.org/abs/2405.15793}

\bibitem{agentless}
X. Xia, Y. Deng, S. Dunn, L. Zhang.
\newblock Agentless: Demystifying LLM-based Software Engineering Agents.
\newblock arXiv:2407.01489, 2024.
\newblock \url{https://arxiv.org/abs/2407.01489}

\bibitem{reflexion}
N. Shinn, F. Cassano, E. Berman, A. Gopinath, K. Narasimhan, S. Yao.
\newblock Reflexion: Language Agents with Verbal Reinforcement Learning.
\newblock arXiv:2303.11366, 2023.
\newblock \url{https://arxiv.org/abs/2303.11366}

\bibitem{sweci}
J. Chen, X. Xu, H. Wei, C. Chen, B. Zhao.
\newblock SWE-CI: Evaluating Agent Capabilities in Maintaining Codebases via Continuous Integration.
\newblock arXiv:2603.03823, 2026.
\newblock \url{https://arxiv.org/abs/2603.03823}

\bibitem{metamon}
H. Lee, G. An, S. Yoo.
\newblock METAMON: Finding Inconsistencies between Program Documentation and Behavior using Metamorphic LLM Queries.
\newblock arXiv:2502.02794, 2025.
\newblock \url{https://arxiv.org/abs/2502.02794}

\bibitem{treesitter}
Tree-sitter.
\newblock Tree-sitter: An Incremental Parsing System for Programming Tools.
\newblock \url{https://github.com/tree-sitter/tree-sitter}

\end{thebibliography}

\end{document}
